\documentclass[10pt, letterpaper]{article}
\usepackage[ignoreheadfoot,top=2cm,bottom=2cm,left=2cm,right=2cm,footskip=0.8cm]{geometry}
\usepackage{fontawesome5}
\usepackage{hyperref}
\hypersetup{colorlinks=true,urlcolor=black}
\usepackage[french]{babel}
\usepackage[T1]{fontenc}
\usepackage{lmodern}
\pagestyle{empty}
\setlength{\parindent}{0pt}
\setlength{\parskip}{0.5em}
\begin{document}
\begin{center}
\textbf{\LARGE Objet : Candidature -- Quantitative Risk Analyst (Financial Services)}\\[0.2cm]
\textbf{\large Charly-Romy TANGA}\\[0.2cm]
\href{mailto:charlyromytanga.cytech@gmail.com}{\faEnvelope\ charlyromytanga.cytech@gmail.com} \quad$\cdot$\quad
\faPhone\ +33 6 15 42 25 74 \quad$\cdot$\quad
\href{https://www.linkedin.com/in/charly-romy-tanga}{\faLinkedin\ charly-romy.tanga} \quad$\cdot$\quad
\href{https://github.com/charlyromytanga}{\faGithub\ charlyromytanga}
\end{center}
Madame, Monsieur,

Je souhaite orienter ma carrière vers le quantitative risk analysis appliqué aux services financiers, avec un fort intérêt pour le risque réglementaire, les métriques de risque fonds et le reporting structuré.

Mon parcours d'ingénieur m'a permis de développer des compétences solides en modélisation statistique, simulation Monte-Carlo et automatisation de pipelines de données en Python/SQL.

Je suis particulièrement attiré par le contexte luxembourgeois, où la stabilité, la conformité et la qualité de la documentation sont des priorités majeures.

Je souhaite contribuer à des missions de risk analytics, de validation de modèles et de transformation data au sein de votre organisation.

Cordialement,\\
Charly-Romy TANGA
\end{document}
